\section{Adding zero-knowledge}
\label{s:zk}


\subsection{Randomizing $w_M(\vec\alpha)$.}
\label{s:randomizing:w_M}

Recall that in the end of Phase \ref{i:main:phase:2} of Protocol \ref{prot:main}, the prover reveals the values
$w_A(\vec\alpha)$, $w_B(\vec\alpha)$, $w_C(\vec\alpha)$, which are subject to 
\begin{equation}
\label{e:zerocheck:last:check}
q_{m}(\alpha_m) = (w_A(\vec\alpha)\cdot w_B(\vec\alpha) - w_C(\vec \alpha))\cdot eq(\vec\alpha,\vec\rho),
\end{equation}
where $\alpha\in F^m$ are the random challenges collected over the rounds of the sumcheck.
In this section we discuss two different strategies to prevent these values from leaking private witness data.
%Both randomize either $w_1$ or $w_2$ ``on-chip'', by 
Both extend the witness (either $w_1$ or $w_2$) by random values, and add trivial gates to assure their propagation through the R1CS matrices. 



\subsubsection{Randomization by extension field elements.}

% This is the canonical solution, where the openings over the extension field $F$ are randomized by random witnesses over $F$. 
Although we do not use this technique in our current implementation, we keep this section for an optional shift to a the second stage over $F$. 
%This breaks with our current design decision keeping $w_2$ also over $\mathbb F_q$, so the solution would be only interesting if we decide to revert to $w_2$ over the extension field $F$.   
In that case, randomizing the openings by additional random witnesses over $F$ is the natural way to go.

First of all, note that it is sufficient to randomize $w_A(\vec\alpha)$ and $w_B(\vec\alpha)$; the value of $w_C(\vec\alpha)$ is then uniquely determined by \eqref{e:zerocheck:last:check}, as long as $eq(\vec\alpha, \vec\rho)\neq 0$. 
We allocate two private extension field variables $R_A, R_B$ in the second stage R1CS, and set their values to $r_A, r_B \sample F$. 
Their gates are 
\[
R_A \cdot 1 = R_A, \quad 1 \cdot R_B = R_B, 
%\quad 1 \cdot R_C = R_C,
\]
located at arbitrary row indices $\vec x_A, \vec x_B \in H_m$.
Then 
\begin{align}
\label{e:wA:randomized:extfield}
w_A(\vec Y) &= w_A'(\vec Y)  + r_A \cdot eq(\vec x_A,\vec Y),
\\
\label{e:wB:randomized:extfield}
w_B(\vec Y) &= w_B'(\vec Y) + r_B \cdot eq(\vec x_B, \vec Y) 
\end{align}
with $w_A'$, $w_B'$, not depending on the two random witnesses $r_A$, $r_B$.
 
% These two randomnesses are sufficient to hide the other witnesses, for most of the $\vec\alpha\in F^m$.
% $\vec x_1\in H_m$ corresponds to the row index of the gate $R_1\cdot 1 = R_1$ in $A_2$.

\begin{lem}
\label{lem:ext:field:randomization}
For any $\vec\rho\in F^m$.
Except for a fraction of $\varepsilon \leq \frac{3 m}{|F|}$ points $\vec\alpha\in F^m$, 
the distribution of $(w_A(\vec\alpha), w_B(\vec\alpha), w_C(\vec\alpha))$ is uniform on the set
\[
\left\{
(v_A, v_B, v_C) \in F^3 
\::\:
q_m(\alpha_m) = (v_A\cdot v_B - v_C)\cdot eq(\vec\alpha,\vec\rho)
\right\}.
\]
\end{lem}
\begin{proof}
The exceptional set is the union of the roots of the multilinears $eq(\vec x_A, \:.\:)$, $eq(\vec x_B, \:.\:)$  and $eq(\:.\:,\vec\rho)$.
For all other $\vec\alpha\in F^m$, the vector $(w_A(\vec\alpha), w_B(\vec\alpha))$ is uniformly distributed over $F^2$, since the coefficients of $r_A$ and $r_B$ do not vanish in \eqref{e:wA:randomized:extfield} and \eqref{e:wB:randomized:extfield}, and they uniquely determine $w_C(\vec\alpha)$ via the claim \eqref{e:zerocheck:last:check}.
$\Box$
\end{proof}

% We note that this already diverges from the univariate case, where local randomization of the witness propagates to \textit{all} points outside the witness domain. 

The bound in Lemma \ref{lem:ext:field:randomization} is essentially tight, as multivariate Lagrangians have exactly a fraction of $m/|F|$ zeros in $F^m$.
In other words, local randomization of the witness propagates only to \textit{almost all} points outside the witness domain.\footnotemark
\footnotetext{%
This diverges from univariate IOPs, where local randomization propagates to \textit{all} points outside the witness domain. 
}
% 
However, the fraction of exceptional points is of order $1/|F|$, and excluding them does not conflict soundness: 
Since $\vec x_A$ and $\vec x_B \in H_m$, we simply exclude the set $\{0,1\}$ from the sampling domain of the zerocheck randomnesses $\alpha_1,\ldots,\alpha_m$, which comes with a negligble increase of the round-by-round soundness errors.

For completeness we also cite an extension to multiple evaluation points:
\begin{lem}
\label{lem:multipoint:eval}
Choose $\Omega =\{\vec x_1, \ldots, \vec x_k\}$ a set of $k\geq 1$ distinct points on the Boolean hypercube $H_n$. 
Then, except with probability of at most $\varepsilon = \frac{k\cdot n}{|F|}$ in $\vec \alpha_1, \ldots, \vec \alpha_k \in F^n$, the multi-point evaluation matrix
\begin{equation*}
E_{\vec\alpha_1,\ldots,\vec\alpha_k} =
\begin{pmatrix}
eq(\vec x_1,\vec \alpha_1) & eq(\vec x_1, \vec\alpha_2) & \ldots & eq(\vec x_1, \vec\alpha_k)
\\
eq(\vec x_2,\vec \alpha_1) & eq(\vec x_2, \vec\alpha_2) & \ldots & eq(\vec x_2, \vec\alpha_k)
\\
\vdots & \vdots & &\vdots
\\
eq(\vec x_k,\vec \alpha_1) & eq(\vec x_k, \vec\alpha_2) & \ldots & eq(\vec x_k, \vec\alpha_k)
\end{pmatrix}
\end{equation*}
is of full rank.
\end{lem}

The lemma follows from that the determinant is a non-zero multilinear in the coordinates of $\vec\alpha_1, \ldots, \vec\alpha_k$.
Since it is not important to us, we leave a detailed proof to the reader.
% \begin{proof}
% The determinant of the matrix equals
% \[
% \sum_{\sigma\in S_k} eq(\vec x_{\sigma(1)}, \vec\alpha_1) \cdot eq(\vec x_{\sigma(2)}, \vec\alpha_2) \cdot \ldots \cdot eq(\vec x_{\sigma(k)},\vec\alpha_k),
% \]
% where $\sigma$ ranges over all permutations of $\{1,2,\ldots, k\}$, and therefore is a multilinear polynomial in the coordinates of $\vec\alpha_1$, \ldots, $\vec\alpha_k$.
% By Schwartz-Zippel, it can be zero for at most $k\cdot n$ points, unless it is identical zero.
% To exclude the latter, observe that each product in the determinant can be written as 
% \[
% L_{\vec x_{\sigma(1)}, \vec x_{\sigma(2)},\ldots, \vec x_{\sigma(k)}}(\vec \alpha_1,\ldots,\vec\alpha_k),
% \]
% the Lagrangian at $(\vec x_{\sigma(1)}, \vec x_{\sigma(2)},\ldots, \vec x_{\sigma(k)})$ in the $k\cdot n$-dimensional hypercube.
% Since the points $\vec x_1,\ldots, \vec x_k$ are distinct, so are the points in 
% \[
% \{(\vec x_{\sigma(1)}, \vec x_{\sigma(2)},\ldots, \vec x_{\sigma(k)}) \::\: \sigma\in S_k \},
% \]
% and we conclude from linear independence of Lagrangians that their sum cannot be zero.
% $\Box$
% \end{proof}




\subsubsection{Randomization by basefield elements.}
%In the light of Lemma \ref{lem:multipoint:eval} 
% As we would like to keep the witness over the basefield $\mathbb F_q$

It is tempting to believe that the univariate strategy from \cite{cryptoeprint:2024/1037}, which randomizes an extension field opening by extension-field-degree many basefield values, carries over to the multilinear world in a straight-forward manner.
However, the best generic result we found is the following one, which has a way too large bound for the exceptional set.

\begin{lem}
\label{lem:Lagrange:random:eval}
Let $F$ be a degree $e$ extension of a prime field $\mathbb F_q$, and $\Omega =\{\vec x_1, \ldots, \vec x_e\}$ a set of $e$ distinct points on the Boolean hypercube $H_m\subset \mathbb F_q^m$. 
Then, except for a fraction of at most $\varepsilon = \frac{m}{q-1}$ points $\vec \alpha \in F^m$, the elements
\begin{equation*}
eq(\vec x_1, \vec \alpha), \ldots, eq(\vec x_e, \vec \alpha) \in F
%= eq(\vec \rho,\vec x_i)    
\end{equation*}
are $\mathbb F_q$-linear independent. 
\end{lem}


The proof of the lemma relies on the well-known Moore determinant as a distinguisher for $\mathbb F_q$-linear independence.
However, we require the property also over $K$ the field of rational functions in several variables, and hence state it in full generality.
%on , which is as follows.
\begin{lem}
\label{lem:goss}
Let $K$ be a (possibly infinite) field of characteristic $q<\infty$, and let $f_1,\ldots,f_k\in K$, $k\geq 1$.
Then $f_1,\ldots, f_k$ are $\mathbb F_q$-linear independent if and only if the determinant of the Moore matrix
\[
% \Delta(f_1,\ldots, f_k) = \left|
M(f_1,\ldots, f_k) =
\begin{pmatrix}
f_1 & f_2 & \ldots & f_k
\\
f_1^q & f_2^q & \ldots & f_k^q
\\
\vdots &\vdots & \ddots & \vdots
\\
f_q^{q^{k-1}} & f_2^{q^{k-1}} & \ldots & f_k^{q^{k-1}}
\end{pmatrix}
% \right| \neq 0
\]
is a non-zero element in $K$. 
\end{lem}

For a proof of Lemma \ref{lem:goss}, we refer the reader to \cite[Lemma 1.3.3]{Goss96}.


\begin{proof}[of Lemma \ref{lem:Lagrange:random:eval}]
Since the points $\vec x_1,\ldots, \vec x_n$ are distinct, the Lagrangians $eq(\vec x_i,\vec Y)$, $i=1,\ldots,e$, are $\mathbb F_q$-linear independent elements of $K= \mathbb F_p(\vec Y)$, the rational function field in $m$ variables. % the coordinates of $\vec Y$.  
By Lemma \ref{lem:goss}, their Moore determinant 
\[
p(\vec Y) = \det M(eq(\vec x_1, \vec Y),\ldots, eq(\vec x_e, \vec Y))\neq 0,
\]
as an element of $K$.
Moreover, since all matrix entries are polynomials, we conclude that the Moore determinant
is a non-zero \textit{polynomial}. 

% The degree in the first coordinate $y_1$ is at most $\frac{q^e - 1}{q-1} \approx q^{e-1}$ 

% Then
% \[
% \prob{y_1\in F \::\: p(y_1, \:.\:) = 0 } \leq 1/q.
% \]
% Taking independence of 

% (All entries of the Moore matrix are polynomials.)
%
% Let us determine its total degree. 
The $(i,j)$-th entry in the Moore matrix is
% \[
% L_{x_j}(Y_1,\ldots, Y_n)^{p^{i}} = L_{x_j}(Y_1^{p^i},\ldots, Y_n^{p^i}),
% \]
has total degree at most $m\cdot q^{i-1}$, and the total degree the determinant is at most
\begin{align*}
m + m q + \ldots + m  q^{e-1} = m \cdot \frac{q^e - 1}{q - 1}.
\end{align*}
By Schwartz-Zippel, the fraction of $\vec\alpha\in F^m$ for which the determinant evaluates zero
is at most
\[
\frac{m}{q-1} \cdot \frac{q^e -1}{q^e} < \frac{m}{q-1}.
\]
For all other $\vec\alpha$ the determinant is non-zero, and another application of Lemma \ref{lem:goss} yields that $eq(\vec x_1,\vec\alpha),\ldots, eq(\vec x_e, \vec \alpha)$ are $\mathbb F_p$-linear independent elements of $F$.
$\Box$
\end{proof}

It is hard to tell whether the Schwartz--Zippel bound in Lemma \ref{lem:Lagrange:random:eval} is tight. 
A high polynomial degree is not a guarantee for a large set of zeros. 
% For example, 
% \[
% eq(\vec x_i, \vec Y)^{q^{i-1}} = eq(x_{i,1},\ldots, x_{i,m}, Y_1^{q^{i-1}},\ldots, Y_m^{q^{i-1}})
% \]
% is of total degree $m\cdot q^{i-1}$. 
% Its zeros are the roots of $Y_j^{q^{i-1}} - x_{i,j}$, but in $F$ there is only one solution, $Y_j = x_{i,j}$, since is no $q^{i-1}$-th root of unity, except $1$.\footnotemark
% \footnotetext{%
% This follows from $(q^{i-1},q^{e} - 1) = 1$. 
% }
% Thus the polynomials has the same number of roots as $eq(\vec x_i, \vec Y)$. 
We tried shortly with GPT-5.6, but did not achieve significant improvements. 
However, it proposed the following work around, which explicitly uses the multiplicative structure of the Lagrangians. 

\paragraph{Blowing up soundness.}
Instead of a single set $\Omega$ of size $e$, we choose several sets
\[
\Omega_1, \ldots \Omega_s \subset H_m,
\]
in a very particular manner, so that their exceptional sets are statistically independent.
(The latter will be possible because of the tensor structure of the Lagrangians.)
Each $\Omega_j$ is a set of $e$ points contained in a sub-cube of small dimension 
\[
t = \ceil{\log_2 e},
\]
with its free part touching the variables $Y_{i_j(1)}, \ldots, Y_{i_j(t)}$ only, meaning that its points do not vary outside these $t$ coordinates.  
%We call the variables $Y_{i_j(1)}, \ldots, Y_{i_j(s)}$ \textit{free variables}, and the remaining $m-s$ variables the \textit{fixed variables} of $\Omega_j$.
We further require that the free variables are disjoint between the sets, i.e.
\[
\big\{Y_{i_j(1)},\ldots, Y_{i_j(t)} \big\} \cap \big\{Y_{i_{j'}(1)},\ldots, Y_{i_{j'}(t)} \big\} = \varnothing,
\]
whenever $j\neq j'$.
The crux is that for each such set $\Omega_j$, we can drop its fixed variables in the proof of Lemma \ref{lem:Lagrange:random:eval}, as long as their $eq$-terms are not zero, since a common non-zero factor does not change $\mathbb F_q$-linear independence. 
The exceptional set for $\Omega_j$ is only in the variables $Y_{i_j(1)}, \ldots, Y_{i_j(t)}$, independent of the other variables, and we arrive at the following statement.

\begin{lem}
\label{lem:zk:basefield:randomization}
With $\Omega_1,\ldots, \Omega_s \subset H_{m}$ as above.
Then except for an $\varepsilon$-fraction of points $\vec\alpha\in F^m$, where
\begin{equation}
\label{e:basefield:randomization:error}
\varepsilon < s\cdot \frac{m-t}{|F|} + \left(\frac{t}{q-1}\right)^s,
\end{equation}
at least one of the sets $\Omega = \Omega_j$ is such that its values $eq(\vec x, \vec\alpha)$, $\vec x\in \Omega$, are $\mathbb F_q$-linearly independent.
\end{lem}
\begin{proof}
The first term bounds the the probability that $\vec\alpha\in F^m$ hits one of the $m-t$ fixed variables of $\Omega_j$, times $s$ for the union bound over all sets.
For all other $\vec\alpha$, the exceptional sets in Lemma \ref{lem:Lagrange:random:eval} are statistically independent, and of relative size $t/(q-1)$ each.
$\Box$
\end{proof}

% \begin{rem}
% \label{rem:propagation:constraints}
In our protocol we will use $s = e$ sets $\Omega_1,\ldots, \Omega_e$, to make the second term in \eqref{e:basefield:randomization:error} of order $O(1/|F|)$. 
That is,  we require overall $e^2$ random witnesses over $\mathbb F_q$ to randomize $w_A(\alpha)$, and another $e^2$ random witnesses for $w_B(\alpha)$. 
The exceptional set for which all $e$ Moore determinants vanish, will be filtered out by the prover.
See the protocol description in Section \ref{s:complete protocol}.
%They can be placed at arbitrary positions in the witness, yet their constraints (as before of the form $R_j \cdot 1 = R_j$, or $1\cdot R_j = R_j$) have to respect the positioning described above. 
% \end{rem}
%The main disadvantage is that we cannot place them at arbitrary positions. 

% d

\subsection{Zero-check masking}
\label{s:zk:zerocheck}

The zero-check, Phase \ref{i:main:phase:2} in Protocol \ref{prot:main}, is turned into zero-knowledge by using the sumcheck randomization from Libra \cite{cryptoeprint:2019/317}.
The prover commits to the univariate polynomials 
\begin{equation}
\label{e:zk:zerocheck:randomizers}
h_0\sample F, \quad h_1(X),\ldots, h_m(X)\sample F[X]^{<d},
\end{equation}
where $d=3$ is the degree of the sumcheck \eqref{e:R1CS:zerocheck}.
These polynomials define the sparse mask
\begin{equation}
\label{e:zk:zerocheck:sparse:mask}
h(X_1,\ldots, X_m) = h_0 + X_1\cdot h_1(X_1) + \ldots + X_m\cdot h_n(X_m),
\end{equation}
which is overlaid with the hypercube sum \eqref{e:R1CS:zerocheck}. 
The prover claims the sum $c' = \sum_{\vec x \in H_m} h(\vec x)$, the computation of which is almost for free via the succinct representation
\begin{align*}
\frac{1}{2^m}\sum_{\vec x\in H_m} h(\vec x) 
% = 2^{m} h_0 + 2^{m-1} h_1(1) + 2\cdot \sum_{\vec x\in H_{m-1}} X_2 h_2(x_2) + \ldots + x_m\cdot h_m(x_m) 
% \\
%  = 2^{m} h_0 + 2^{m-1} h_1(1) + 2\cdot (2^{m-2} h_2(1) + 2 \cdot \sum_{\vec x H_{m-2}} ...)
%  \\
=  h_0 + \frac{1}{2}\left(
h_1(1) + \ldots + h_m(1)
\right).
\end{align*}
%It receives a random $\alpha\sample F$ from the verifier, and 
Then both run the sumcheck on the combined claim
\begin{equation}
\label{e:zk:zerocheck:combined:claim}
c_0 = \alpha \cdot c + c' = \sum_{\vec x\in H_n} \alpha\cdot G(\vec x) + h,
\end{equation}
where $\alpha\sample F$ is random, and  $c = \sum_{x\in H_n} G(x_1,\ldots, x_n)$ is the target claim. 
(In our case $G$ is be the multivariate zero-check term, and $c=0$.)

%
The sparse mask carries sufficient entropy to perfectly hide the witness in the sumcheck polynomials
\begin{equation*}
\hat q_k(X) = \sum_{\vec x\in H_{n-k}} \alpha\cdot G(\alpha_1,\ldots, \alpha_{k-1},X,\vec x) +  h(\alpha_1,\ldots, \alpha_{k-1},X,\vec x),
\end{equation*}
for $k=1,\ldots,n$, as the following lemma shows.

\begin{lem}
For any $\alpha\in F $, $(\alpha_1,\ldots, \alpha_n)\in F^n$. 
The distribution of the sumcheck prover messages $c_0, \hat q_1(X),\ldots, \hat q_n(X)$ is uniform on the set
\[
\mathcal R = \left\{
\begin{matrix}
c_0\in F, \hat q_1,\ldots,\hat q_n\in F[X]^{\leq d}
\end{matrix}
\::\:
\begin{matrix}
c_0 = \hat q_1(0) + \hat q_1(1),
\\
\hat q_{k}(\alpha_{k}) = \hat q_{k+1}(0) + \hat q_{k+1}(1),
\\
k=1,\ldots n-1
\end{matrix}
\right\},
\]
independent of $G$.
\end{lem}

For the sake of completeness, we give a proof the lemma.
\begin{proof}
Fix $G$, $\alpha\in F$ and $(\alpha_1,\ldots,\alpha_n)\in F^n$, and choose any set $S\subset F$ of $d$ distinct points, containing the element $1$, but not $0$.
Using the consistency equations it is easily seen that the linear mapping, which sends $(c_0,\hat q_1(X),\ldots, \hat q_n(X))$ to the values   
\begin{equation}
\label{e:zk:zerocheck:polys:by:value}
(\vec v_1,\ldots, \vec v_n) = (\hat q_1|_{S\cup\{0\}},\hat q_1|_S, \ldots, \hat q_n|_S) 
\end{equation}
is a bijection between $\mathcal R$ and $F^{dn + 1}$. 
We leave this to the reader.\footnotemark
\footnotetext{%
The values $\vec v_1$ uniquely determine the polynomial $\hat q_1(X)\in F[X]^{\leq d}$ by degree, and with it $c_0 = \hat q_1(0) + \hat q_1 (1)$.
Likewise, $\hat q_1(\alpha_1)$ together with the values $\vec v_2$ uniquely determine $\hat q_2(X)\in F[X]^{\leq d}$ satisfying $\hat q_1(\alpha_1) = \hat q_2(0) + \hat q_2(1)$, since the letter determines the value of $\hat q_2(X)$ at the additional point $0$.
And so on.
}
To prove the lemma, it suffices to show that the values \eqref{e:zk:zerocheck:polys:by:value} are in one-to-one correspondence with $h_0\in F, h_1,\ldots, h_n\in F[X]^{<d}$.
But this is evident from the concrete form
\begin{multline*}
\hat q_{k}(X) = \sum_{\vec x \in H_{n-k}} \alpha \cdot G(\alpha_1,\ldots, \alpha_{k-1}, X,\vec x) 
\\
+ 2^{n-k} (h_0 + \alpha_1 h_1(\alpha_1) +\ldots +\alpha_{k-1} h_{k-1}(\alpha_{k-1})+ X h_{k}(X)) 
\\+ \sum_{\vec x \in H_{n-k}} x_{k+1} h_{k+1}(x_{k+1}) + \ldots + x_n 
h_n(x_n),
\end{multline*}
in which the other $h_0,\ldots, h_{k-1}, h_{k+1},\ldots, h_n$ impact only the constant part of the polynomial.
(The representation partially collapses in the two edge cases.
For $k=1$ there are no terms except $h_0$ preceding $X h_1(X)$, and for $k=n$, there are no terms following $X h_n(X)$.)
$\Box$
\end{proof}


The sumcheck reduces the claim \eqref{e:zk:zerocheck:combined:claim} to that 
\begin{align*}
\hat q_m(\alpha_m) = G(\alpha_1,\ldots, \alpha_m) + h_1(\alpha_1) + \ldots + h_m(\alpha_m),
\end{align*}
where
\begin{align*}
G(\vec\alpha) =  (w_A(\vec\alpha)\cdot w_B(\vec\alpha) - w_C(\vec\alpha))\cdot eq(\vec\alpha,\vec\rho),    
\end{align*}
for which the prover sends the values $w_A(\vec\alpha)$, $w_B(\vec\alpha)$, $w_C(\vec\alpha)$. 
The complementary term 
\begin{equation}
\label{e:zk:zerocheck:mask:opening}
h_1(\alpha_1) + \ldots + h_m(\alpha_m) = \hat q_m(\alpha_m) - G(\vec\alpha)
\end{equation}
can be proven in zero-knowledge by a separate protocol.

\begin{rem}
In our current implementation, we append the $dm+1$ coefficients\footnotemark for $h_0,h_1,\ldots, h_m$ to the first stage witness $\vec w_1$, and use zk-WHIR with the additional constraint
\begin{equation}
\hat q_m(\alpha_m) - G(\vec\alpha) = \sum_{\vec x\in H_{m_1}} I_\alpha(\vec x)\cdot  w_1(\vec x),
\end{equation}
where $I_\alpha$ is the inner product kernel for $h_1(\alpha_1) + \ldots + h_m(\alpha_m)$.
This requires the witness to be sufficiently below a power two. 
See the complete protocol in Section \ref{s:complete protocol}.
\footnotetext{%
Alternatively, one may directly use the representation by their values \eqref{e:zk:zerocheck:polys:by:value}.} 
\end{rem}

Let us quickly sketch the transcript simulator for the zero-check.
The prover chooses arbitrary random polynomials $w_1(X_1,\ldots, X_{m_1})$  and $w_2(X_1,\ldots, X_{m_2})$, with their image polynomials $w_A, w_B, w_C$ not satisfying the zero-check hypercube sum for random $\vec\rho\sample F$.
Then it samples $h_0,h_1,\ldots, h_n$ and $\lambda\neq 0$ and computes the hypercube sum
\begin{equation}
\label{e:sumcheck:randomized:G}
c_0 = \sum G + \lambda h.
\end{equation}
It then runs the sumcheck protocol on the valid claim \eqref{e:sumcheck:randomized:G}.
The (wrong) claim $c'$ is chosen adaptively so that $\alpha \cdot 0 + c' = c_0$. 
For details, see Section \ref{}.

% For on-domain randomization we rely on the following Lemma. 
% \ulrich{We need an adaption to multivariates of bounded individual degree, if we reveal cross-terms}
% \begin{lem}
% Fix any $k$ distinct points $\vec x_1,\ldots, \vec x_k\in H_n$. 
% Then 
% \[
% \Pr_{\vec y_1,\ldots, \vec y_k\sample F^n}\left[
% \det \begin{pmatrix}
% eq(\vec y_i,\vec x_j)
% \end{pmatrix}
% = 0
% \right] 
% \leq \frac{k\cdot n}{|F|}.
% \]
% \end{lem}
% \begin{proof}
% The determinant is multilinear in the coordinates of $\vec y_1, \ldots, \vec y_k$.
% $\Box$
% \end{proof}


\subsection{Overlaid masks instead of zk-encoding.}
\label{s:masked:commitments}

To avoid witness information leakage from commitment openings, we use a slightly different approach than \cite{cryptoeprint:2026/391}.
Instead of committing to a zero-knowledge encoding of witness,
we overlay a witness polynomial with a sparse random mask without  degree increase:
Given witness $p \in \mathscr P_{m}$ and random $\msk\sample F[X]^{< q_i}$, for some small $q_i < 2^m$, we commit to
\begin{equation}
\hat p = p + \msk, %\in \mathscr P_m,
\end{equation}
which is again from $\mathscr P_m$.
(This secures the witness for up to $q_i$ openings.)
Obviously this does not bind the witness, unless we also commit to $\msk$, which is required for ZO0K anyways (see Section \ref{s:zook}).

For practical reasons, we \textit{jointly} commit all\footnotemark masks required by the protocol in one shot.
\footnotetext{%
These are the masks for the two witness polynomials $w_1$ and $w_2$, and the further masks $\msk_{1},\ldots, \msk_{r}$, $h_1,\ldots, h_r$ required for ZO0K.
}%
We concatenate their coefficient vectors, append $q_{zk}$ additional random coefficients, and encode the resulting vector as a single word from 
\begin{equation}
\label{e:C:zk}
C_{zk} = \RS[F,D_{zk}, d_{zk}].
\end{equation}
Here the degree bound $d_{zk}$ is at least $q_{zk}$ larger than the total of all mask lengths, and can be chosen again a power of two, if this helps encoding performance.
See Section \ref{s:zook} for details.
(The choice of $q_{zk}$ depends on the number of openings in the final proof of proximity, Protocol \ref{prot:zook:final:opening}.)

The rationale behind masked commitments is the following.
Adding the small mask $\msk$ in coefficient representation is a local modification of the witness, and the 
subsequent encoding of $\hat p$ is as costly as that of $p$. 
In contrast, even a tiny degree increase of the witness polynomial would lead to an $O(m)$ overhead in the encoding step, where $m$ is the degree of $p$.




\subsection{ZO0K}
\label{s:zook}


We describe an adaptation of ZO0K \cite{cryptoeprint:2026/391} to the masked commitments as described in the previous section.
The changes to the original protocol are minimal, and based on the fact that we can now also view masked polynomials as multilinears in the sumcheck claim.
See Remark \ref{rem:zook:pi:instead:pi:hat} for a one-to-one variant of \cite{cryptoeprint:2026/391}.


% \subsubsection{ZO0K.}
%ZO0K instead uses throughout sparse masks for randomization, and is therefore particularly prover friendly. 
The sumcheck rounds are randomized using the sparse Libra mask explained in Section \ref{s:zk:zerocheck}, and information leakage via commitment openings is mitigated through the commitment masks as described above.
In this multi-polynomial regime, ZO0K considers WHIR as a sumcheck prover, which gradually reduces composite claims of the form 
\begin{align}
\label{e:zook:claim}
c_{i}  &= \left\langle \hat p_{i}, K^{(i)} \right\rangle_{H_{k_{i}}} 
+ \sum_{j=0}^{i} \Gamma_j^{(i)}(\msk_j) + \sum_{j=1}^{i} \Lambda_{j}^{(i)} (h_{j}),
\end{align}
involving a costly part, the inner product $\left\langle \hat p_{i}, K^{(i)} \right\rangle$ of a large (masked) witness polynomial $\hat p_i$ with a public kernel $K^{(i)}$, and a cheap part, certain public linear functionals $\Lambda_j^{(i)}$, $\Gamma_j^{(i)}$ on the (small) masks committed by the prover so far:\footnotemark
\footnotetext{%
The $\Lambda_j^{(i)}(h_j)$ and $\Gamma_j^{(i)}(\msk_i)$ can be again represented as inner products. 
However, because of their fundamentally different treatment in the protocol, we write them as linear functionals.
}
The sparse sumcheck masks $h_1,\ldots, h_i$ and the witness masks $\msk_0,\ldots,\msk_i$.
The cheap part is not further reduced by the sumcheck logic; it is treated as constant function 
% \[
% \frac{1}{2^{k_i}}\cdot \sum_{j=1}^{i-1} \Lambda_{j}^{(i)} (h_{j}) + \Gamma_j^{(i)}(\msk_j)
% \]
on the hypercube $H_{k_i}$, which is passed through the rounds in a weighted form.
%leading to an simple update law for the linear functionals. 
The costly part is reduced in the usual sense, by folding. 
However, not before a fresh sumcheck mask $h_{i+1}$ enters the sum.
In this manner, the number of masks in the claim \eqref{e:zook:claim} grow from round to round, but the costly part shrinks.

After the final round, the remaining claim, a linear functional in $\hat p_r$, $h_1,\ldots, h_r$, $\msk_0,\ldots, \msk_{r}$, is proven in zero-knowledge by an elementary, non-succinct proof of proximity, Protocol \ref{prot:zook:final:opening} below.




% Recall that WHIR in full generality is a sumcheck prover, and works for arbitrary algebraic expressions $G(f_1,\ldots, f_m)$ in several multilinear polynomials $f_1$, \ldots $f_m$. 
% Every now and then it accompanies the sumcheck with new proposals for the multilinears in the specialized claim, and the claim of proximity to a sumcheck consistent ensemble of polynomials is reduced to a fresh claim on the new ensemble of smaller codewords.
% However there is no strict need to provide new proposals for every multilinear in the claim. 
% If a multilinear is small, e.g. depending only one or two variables or has a succinct representation, then the verifier can compute its specialization on its own, without sacrificing efficiency.

% This is the view point we take on ZO0K:
% In each round the prover replaces the (one and only) large polynomial in the claim by a smaller one. 
% On the other hand it adds few more small commitments required for zero-knowledge.
% The number of committed words grows, but the total size of the multilinears becomes smaller and smaller.


% (Here, we again use the same notation for the composed polynomial and its interleaved representation.)
% This is a linear constraint on the committed codewords $\hat p_0$ and $\hat\msk_0$.
 %(Their final specialized claim will be proven on their initial unfolded form.)
% Once the polynomial becomes 

\subsubsection{The required masks.}
As pointed out in Section \ref{s:masked:commitments}, we assume that the prover has committed beforehand to all the masks, as a single word from $C_\text{zk}$.
Concretely, each sumcheck mask $h_i$ is committed as coefficient vector 
\begin{equation}
\label{e:zook:sumcheck:mask:coeffs}
h_{i,0}\| h_{i,1} \| \ldots\| h_{i,r_i} \in F^{2 r_i + 1}
\end{equation}
from $h_{i,0}\sample F$, and $h_{i,j}\sample F[X]^{<2}$, appended by the coefficients of the commitment masks \begin{align}
\label{e:zook:commitment:mask:coeffs}
msk_i \in F^{q_i},  \quad i=0,\ldots,r,
\end{align}
from $\msk_i \sample F[X]^{< q_i}$, where $msk_0$ is taken from the subfield $\mathbb F_q$, assuming that the first witness polynomial $p_0$ is over $\mathbb F_q$.
(In our application, our $p_0$ is the composite witness polynomial, both of which have their own basefield mask, and $msk_0 = msk_{0,1} \| msk_{0,2}$.)
The number of coefficients is chosen according to how often the commitment of $p_i + \msk_i$ is queried in the course of WHIR:
\begin{align}
q_0 &= s_{1}, 
\\
q_i &= s_{i+1} + 1 \text{ for } 1\leq i\leq r-1,
\\\intertext{assuming a single DEEP query in each of the rounds, and}
q_r &= s_\text{final},
\end{align}
the number of queries on $\hat p_r$ in the final opening protocol, Protocol \ref{prot:zook:final:opening}.
The concatenated coefficient vectors \eqref{e:zook:sumcheck:mask:coeffs} and \eqref{e:zook:commitment:mask:coeffs} are encoded in zero-knowlege manner, by appending to the message another $q_{zk} = s_\text{zk} + 1$ random coefficients
\begin{equation}
\label{e:zook:mask:blinder:coeffs}
msk_{zk} \in F^{q_{zk}},
\end{equation}
before the Reed--Solomon encoding. 
For reasons to be explained later, another $e =[F:\mathbb F_q]$ random elements $r_0\sample \mathbb F_q^e$ are appended to it.
Overall, $C_{zk}$ encodes the joint mask vector  
\begin{equation}
\label{e:all:masks:concatenated}
\mathsf m = h_1\|\ldots \|h_r \| msk_0\| \ldots \| msk_r\| msk_{zk} \| r_0
\end{equation}
parsed as coefficient vector of $F^{m_\text{zk}}$,
\begin{equation}
\label{e:Czk:message:length}
m_\text{zk}= r + k_0 - k_r + q_0  + \ldots + q_r + q_{zk} + e,
\end{equation}
as Reed--Solomon codeword over $F$ with evaluation domain $D_{zk}$.


\subsubsection{Initialization.}
The initial claim is on the unmasked polynomial $p_0 = \hat p_0 - \msk_0 \in \mathscr P_{k_1}^{e_1}(\mathbb F_q)$, and of the form
\[
c_0 = \langle \hat p_0 - \msk_0, K \rangle_{H_{k_0}},
\]
where the inner product kernel $K$ is public.
This is of the form \eqref{e:zook:claim} with 
\[
K^{(0)} = K, \quad \Lambda_0^{(0)}(\msk_0) = \langle K, \msk_0\rangle_{H_{k_0}}.
\]
The latter can be easily computed as a linear functional on the coefficient vector $msk_0\in\mathbb F_q^{q_0}$:
The multilinear representation of $\msk_0\in F[X]^{< q_0}$ is non-zero only on the first $q_0$ points of the hypercube and thus
\[
\langle K, \msk_0\rangle_{H_{k_0}} = \sum_{j=0}^{q_0-1} K(j)\cdot msk_{0,j},
\]
in the index notation of the hypercube. 

\subsubsection{Modifications of a WHIR round.}

Let us describe the changes of round $i$ in Protocol \ref{prot:WHIR}, $i=1,\ldots, r$.
At the beginning of the round, the prover has already committed to 
$\hat p_{i-1}$, $\msk_0,\ldots, \msk_{i-1}$, $h_1,\ldots h_{i-1}$, satisfying the linear constraint
\begin{align}
\label{e:zook:prev:claim}
c_{i-1}  &= \left\langle \hat p_{i-1}, K^{(i-1)} \right\rangle_{H_{k_{i-1}}} 
 + \sum_{j=0}^{i-1} \Gamma_j^{(i-1)}(\msk_j)
+ \sum_{j=1}^{i-1} \Lambda_{j}^{(i-1)} (h_{j})
% + \sum_{j=1}^{i-1} \sum_{k=0}^{r_j} \Lambda_{i,j}^{(i-1)} (h_{j,k})
%     + \sum_{j=1}^{i-1} \Gamma_i^{(i-1)}(\msk_j)
\end{align}
We think of the claim as $c_{i-1} = \sum_{\vec x\in H_{k_{i-1}}} G_{i-1}(\vec x)$ for the quadratic expression
\begin{align*}
G_{i-1}(\vec x) = \hat p_{i-1}(\vec x) \cdot K^{(i-1)}(\vec x) - 2^{-k_{i-1}}\cdot B_{i-1},
\end{align*}
where $B_{i-1}$ is the sum of the mask functionals in \eqref{e:zook:prev:claim}.


\paragraph{Step 1: Masking phase.}
Identical to the sumcheck randomization in Section \ref{s:zk:zerocheck}.
The prover takes the masks $h_{i,0}\in F$, $h_{i,1},\ldots,h_{i,r_i} \in F[X]^{<2}$ it committed before running the protocol. 
They define the sumcheck mask
% \ulrich{mayb /}
\begin{equation}
\label{e:zook:sumcheck:mask}
h_i(X_1,\ldots, X_{r_i}) = h_{i,0} + X_1 \cdot h_{i,1}(X_1) + \ldots + X_{r_i} \cdot h_{i,r_{i}}(X_{r_i}),
\end{equation}
in as many variables as the sumcheck phase of the WHIR round consumes, $r_i = k_i - k_{i-1}$.
% in the next $r_i$ variables $X_1,\ldots, X_{r_i}$ processed by the sumcheck.
It claims 
\[
d_i = \sum_{\vec x\in H_{k_{i-1}}} h_i(x_1,\ldots, x_r).
\]
The verifier sends a random $\lambda_i\sample F$, which reduces  \eqref{e:zook:prev:claim} to the combined claim
\begin{align}
\label{e:zook:combined:sumcheck:claim}
\lambda_i\cdot c_{i-1} + d_i = \sum_{\vec x\in H_{k_{i-1}}} 
\left(\lambda_i \cdot G_{i-1}(\vec x) + h_i(x_1,\ldots, x_{r_i})
\right).
\end{align}

\paragraph{Step 2: Sumcheck phase.}
Both prover and verifier proceed with the $r_i$ sumcheck rounds on the combined claim \eqref{e:zook:claim:combined}.
In the course of it, the prover provides the quadratic refinement polynomials
$
\hat q_{i,1}(X), \ldots, \hat q_{i,r_i}(X), 
$ 
the verifier has sent $\vec\beta_i = (\beta_{i,1},\ldots, \beta_{i,r_i})$.
After the last sumcheck round, the claim \eqref{e:zook:combined:sumcheck:claim} is reduced to 
\begin{multline*}
\hat q_{i,r_i}(\beta_{i,r_i}) 
= \sum_{\vec x'\in H_{k_i}} \left(
\lambda_ i \cdot G_{i-1}(
\vec\beta_i,\vec x') + h_i(\vec\beta_{i})
\right)
\\
= \lambda_i\cdot \left\langle \hat p_{i-1}(\vec \beta_i, \:.\:), K^{(i-1)}(\vec\beta_i,\:.\:)\right\rangle_{H_{k_i}} + 2^{k_i - k_{i-1}}\cdot \lambda_i\cdot B_i + 2^{k_i}\cdot h_i(\vec\beta_i)
\end{multline*}
which linear in $\hat p_{i-1}(\vec\beta_{i},\:.\:)$ and the masks, now including $h_i$.

\paragraph{Step 3: Commit.}
The prover takes $\msk_{i}\in F[X]^{<q_i}$, and commits to %the masked polynomial
\[
\hat p_i = \hat p_{i-1}(\vec\beta_{i},\:.\:) + \msk_i(\:.\:)
\]
as a word from $\mathscr P_{k_{i}}^{e_i}$.
This masked commitment is required also in the last round $i=r$. 
% (For an alternative variant see Remark  \ref{rem:zook:pi:instead:pi:hat}.)
The new ensemble of polynomials, $\hat p_i, \msk_i$, and $\msk_0,\ldots, \msk_{i-1}$, $h_1,\ldots, h_i$, is subject to the linear constraint
\begin{multline}
\label{e:zook:specialized:claim}
\hat q_{i,r_i}(\beta_{i,r_i}) = 
\left\langle \hat p_i - \msk_i,  \lambda_i\cdot K^{(i-1)}(\vec\beta_i, \:.\:) \right\rangle_{H_{k_i}} 
\\
+ 2^{-r_i} \cdot \lambda_i \cdot B_{i-1} + 2^{k_i}\cdot h_i(\vec\beta_i),
\end{multline}
which are of the form \eqref{e:zook:claim}.
%except that $\msk_{i-1}$ does not appear in it.

\paragraph{Step 4: Query phase.}
In the query phase, the verifier asks the masked polynomial $\hat p_i$ for the value at the out-of-domain sample $x_{i,0}\sample F$ and uses the values of the previous $\hat p_{i-1}$ to constrain $\hat p_i - \msk_i$ at the points $x_{i,1},\ldots, x_{i,s_i}\sample D_i$.
In terms of the new ensemble, these claims are the linear constraints 
\begin{align}
\label{e:zook:out:domain:constraint}
\sum_{\vec x' \in H_{k_i}} \hat p_i(\vec x') \cdot L_{x_{i,0}}(\vec x') &= v_0,
\\
\label{e:zook:in:domain:constraint}
\sum_{\vec x'\in H_{k_i}} \left(\hat p_i(\vec x') - \msk_{i}(\vec x')\right)\cdot L_{x_{i,k}}(\vec x') &= v_{k},
\end{align}
for $k= 1,\ldots, s_i$, where $L_{x_{i_k}}$ are the univariate Lagrangians of degree $2^{k_{i-1}}$. 
These constraints are again of the form \eqref{e:zook:claim}.
%now with $\msk_{i-1}$ mentioned the first time. 


% The proximate $\hat p_1$ the multilin rep of which is
% \begin{align*}
% \langle \hat p_1,
% K_{0\lambda_1,\ldots, \lambda_{k}}
% \rangle_{H_{n-k}}
% & = q_k(\lambda_k) - \gamma\cdot h(\lambda_1,\ldots,\lambda_{k}) 
% \\
% & \quad\quad + 
% \langle \msk_1 - \msk_{0,\lambda_1,\ldots,\lambda_k},
% K_{0\lambda_1,\ldots, \lambda_{k}}\rangle_{H_{n-k}}
% \\
% \langle \hat p_1,  L_{Q_i}\rangle_{H_{n-k}} &= v_{Q_0}
% \\
% \langle \hat p_1, L_{Q_i}\rangle_{H_{n-k}} &= (\fold_{\lambda_1,\ldots, \lambda_k} \hat f_0) (Q_i).
% \end{align*}




\paragraph{Step 4: Aggregation of claims.}

The verifier samples a random $\gamma_i\sample F$ for combining the above claims \eqref{e:zook:specialized:claim}, \eqref{e:zook:out:domain:constraint} and \eqref{e:zook:in:domain:constraint} on the new ensemble into a single claim of the form \eqref{e:zook:claim}, with the following linear functionals:
\begin{align}
\label{e:zook:update:K}
K^{(i)} &=  \lambda_i\cdot K^{(i-1)}(\vec\beta_i,\:.\:) + \sum_{k=0}^{s_i} \gamma_i^{k+1} \cdot L_{x_{i,k}}(\:.\:),
\\\label{e:zook:update:Gamma}
\Gamma_{j}^{(i)}  &=  2^{r_i}\cdot\lambda_i\cdot \Lambda_j^{(i-1)}, \quad j=1,\ldots,i-1,
\\\label{e:zook:new:Gamma}
% \Gamma_i^{(i)} &= -\Big\langle \:.\:, \lambda_i\cdot K^{(i-1)}(\vec\beta_i,\:.\:) + \sum_{k=1}^{s_i} \gamma_i^{k+1}\cdot L_{x_{i,k}}(\:.\:) \Big\rangle_{H_{k_i}}   
\Gamma_i^{(i)} &= - \langle \:.\:, K^{(i)} - \gamma_i\cdot L_{x_{i,0}} \rangle_{H_{k_i}}
\\\label{e:zook:update:Lambda}
\Lambda_{j}^{(i)}  
&=  2^{r_i}\cdot\lambda_i\cdot \Lambda_j^{(i-1)}, \quad j=1,\ldots, i-1,
\\\label{e:zook:new:Lambda}
\Lambda_i^{(i)}(h_i) & = 2^{-k_i}\cdot h_i(\vec \beta_i).
\end{align}
As before, the inner product \eqref{e:zook:new:Gamma} on $\msk_i$ and the evaluation \eqref{e:zook:new:Lambda} on $h_i$ are inexpensive to compute.


\begin{rem}
\label{rem:zook:pi:instead:pi:hat}
In closer alignment with \cite{cryptoeprint:2026/391}, one may  think of the claims \eqref{e:zook:claim} in terms of the unmasked polynomial $p_i = \hat p_i - \msk_i$, rather than of the masked one, i.e.
\begin{align*}
% \label{e:zook:claim:paper}
c_{i}  &= \left\langle p_{i}, K^{(i)} \right\rangle_{H_{k_{i}}} 
+ \sum_{j=0}^{i} \Gamma_j^{(i)}(\msk_j) + \sum_{j=1}^{i} \Lambda_{j}^{(i)} (h_{j}).
\end{align*}
In Step 3, the prover commits to $\hat p_i = p_{i-1}(\vec\beta_i,\:.\:) + \msk_i(\:.\:)$ instead of $\hat p_i = \hat p_{i-1}(\vec\beta_i,\:.\:) + \msk_i(\:.\:)$, and the consistency relation for the $s_i$ in-domain queries would then be 
\begin{align*}
\sum_{\vec x'\in H_{k_i}} \left(\hat p_i(\vec x') - \msk_{i}(\vec x') + \msk_{i-1}(\vec\beta_i, \:.\:) \right)\cdot L_{x_{i,j}}(\vec x') &= v_{j}, 
\end{align*}
instead of \eqref{e:zook:in:domain:constraint}, leading to a different update rule of $\Gamma_{i-1}^{(i)}$.
\end{rem}



\subsubsection{The final opening protocol.}

The claim after round $r$ is of the form
\begin{equation}
\label{e:zook:claim:r}
c_r  = \langle \hat p_r, K^{(r)} \rangle_{H_{k_r}} 
+ \sum_{j=0}^r  \Gamma_j^{(r)}(\msk_j) + \sum_{j=1}^r \Lambda_j^{(r)}(h_j),
\end{equation}
involving the final $\hat p_r$, and the joint overall mask $\mathsf m$ from \eqref{e:all:masks:concatenated}, comprised of the oracle masks $\msk_0,\ldots,\msk_r$, and the sumcheck blinders $h_1,\ldots, h_r$.
% with $K = K^{(r)}$, $\Lambda_i = \Lambda_i^{(r)}$ and $\Gamma_i = \Gamma_i^{(r)}$.
The claim is linear in its components, and proven via an independent proof of proximity of the committed words from $F^{D_r}$ and $F^{D_\text{zk}}$ to polynomials $\hat p_r(X)\in \mathscr P_{k_r}^{e_r}$ and $\mathsf m(X) \in F[X]^{<m_\text{zk}}$, respectively, satisfying the joint claim \eqref{e:zook:claim:r}.
Such a proof probes the committed words independently, and the soundness errors are blown up by the combinatorial product
\[
\ell(\theta_r) \cdot \ell(\theta_\text{zk})
\]
of the list sizes determined by the number of samples from $D_r$ and $D_\text{zk}$, respectively.

For improving soundness, we once again apply the DEEP principle as in WHIR, and additionally constrain the prover to a single proximate mask by asking for the value 
$
v =\mathsf m(z)
$
at a random out-of-domain point 
\begin{equation}
z\sample F
\end{equation} at commit time of $\mathsf m$. 
(See the full protocol description in Section \ref{s:complete:protocol}.)
We further reuse the same challenge $z$ for proving invariance of the component $msk_0$ under the Frobenius automorphism $\phi$ of $F/\mathbb F_q$. 
This is required for showing the unmasked proximate of $w$ being basefield-valued, see the discussion in Section \ref{s:zk:soundness} below.\footnotemark
\footnotetext{%
There are alternative options to achieve the same goal.
For example, by committing the $\mathbb F_q$-polynomial $\hat\msk_0$ separately. 
}
The prover provides the value of 
\begin{align}
\hat\msk_0(X) = \msk_0(X) + X^{|\msk_0|}\cdot r_0(X)
\end{align}
at $z$, that is the linear functional
\begin{align}
\label{e:zook:final:opening:V:z}
v' = V_z(\mathsf m) &=  \hat\msk_0(z).
\end{align}
Here, $r_0\sample\mathbb F_q^e$ is the component of the overall mask \eqref{e:all:masks:concatenated} devoted for securing $\msk_0$.
If $\phi(v') = V_{\phi(z)}(\mathsf m)$, then with high probability $\hat \msk_0(X) = \phi(\hat\msk_0)(X)$, proving that $msk_0$ (and $r_0$) is over the basefield $\mathbb F_q$.

\begin{protocol}[Final opening]
\label{prot:zook:final:opening}
The prover commits to full entropy masks $H \sample \mathscr P_r$ and $M\sample F^{m_{zk}}$, as words from $C_r$ and $C_{zk}$, respectively.
It parses $M$ as $H_1\|\ldots\|H_r\|$ $M_0 \| \ldots\|M_r\|M_{zk}\|R_0$ according to \eqref{e:all:masks:concatenated}, and claims
\begin{align}
\label{e:zook:claim:H}
d  &= \langle  H, K^{(r)} \rangle_{H_{k_r}}
+ \sum_{j=0}^r \Gamma_j^{(r)}(M_j) + \sum_{j=1}^r \Lambda_j^{(r)}(H_j),
\end{align}
and
\begin{align}
v_{M} = M(z), 
\quad
v_{M}' = V_{z}(M), 
\quad
v_M'' = V_{\phi(z)}(M),
\end{align}
where $V_z$ is as in \eqref{e:zook:final:opening:V:z}.
The verifier replies with $\lambda\sample F$.
The prover reveals the combined messages $H' = \lambda\cdot \hat p_r + H$, $M' = \lambda\cdot \mathsf m + M$, the latter parsed as $H_1'\|\ldots\|H_r'\|M_0'\|\ldots\|M_r'\|$ $M_{zk}'\|R_0'$.

The verifier queries the commitments of $\hat p_r$ and $\mathsf m$ independently at random points $x_{1}, \ldots x_{s}\sample D_r$ and $y_1,\ldots, y_{s_\text{zk}}\sample D_{zk}$, respectively, and checks collinearity of the revealed values with those of the encodings of $H'$ and $M'$. 
%the linear combination.
If the collinearity checks hold, and if 
\begin{align}
\label{e:zook:claim:combined}
\lambda \cdot c_r + d  &= \langle  H', K \rangle_{H_{k_r}}
+ \sum_{j=0}^r \Gamma_j(M_j') + \sum_{j=1}^r \Lambda_j(H_j'),
\end{align}
as well as
\begin{align}
\lambda \cdot v + v_M &=  M'(z),
\\
\lambda\cdot v'  + v_M' = V_z(M'),
&\quad \lambda\cdot \phi(v') + v_M'' = V_{\phi(z)}(M'),
\end{align}
then the verifier accepts.
\end{protocol}
